\begin{ouverture}[calcul]{8}{Grandeurs composées et unités}{Partie 2 · Algèbre et modèles}
{Vitesse · Débit · Concentration · Énergie}
\objectif{Calculer une grandeur quotient ou produit.}
\objectif{Convertir les unités de temps, de volume et de vitesse.}
\objectif{Contrôler la cohérence d’une formule et d’un résultat.}
\prerequis{Proportionnalité et conversions simples.}
\end{ouverture}
\pagestyle{otMain}\essentiel\fontsize{11}{15}\selectfont
\phantomsection\label{cours-8-1}
\essSection{1}{Des quotients de grandeurs}
\begin{definition}Une vitesse moyenne est une distance divisée par une durée : $v=d/t$. Un débit est un volume divisé par une durée : $q=V/t$. Une concentration massique est une masse de substance divisée par le volume de solution : $c=m/V$. Une masse volumique est la masse totale d’une matière divisée par son volume : $\rho=m/V$.\end{definition}
\begin{exemple}Parcourir 150 km en 2 h 30 min donne $t=2{,}5$ h, puis $v=150/2{,}5=60$ km/h. Un robinet débitant 12 L en 40 s a un débit de $12/40=0{,}3$ L/s.\end{exemple}
\begin{savoirfaire}[Choisir des unités cohérentes]Pour des km/h, utiliser des km et des heures. Pour des m/s, utiliser des mètres et des secondes. $1$ h $=3600$ s, $1$ min $=60$ s. Attention : 1 h 30 min vaut 1,5 h, pas 1,30 h.\end{savoirfaire}
\begin{attention}La vitesse moyenne d’un aller-retour n’est pas en général la moyenne des deux vitesses. Il faut toujours diviser la distance totale par la durée totale.\end{attention}
\clearpage
\phantomsection\label{cours-8-2}
\essSection{2}{Convertir des grandeurs composées}
\begin{propriete}$1$ km/h $=1000/3600$ m/s $=1/3{,}6$ m/s. Pour passer de km/h à m/s, diviser par 3,6 ; dans l’autre sens, multiplier par 3,6. Cette règle se retrouve avec les conversions de distance et de temps.\end{propriete}
\begin{propriete}$1$ L $=1$ dm$^3$, $1$ m$^3=1000$ L et $1$ mL $=1$ cm$^3$. Par exemple $0{,}02$ m$^3$/s $=20$ L/s $=1200$ L/min : on convertit séparément volume et temps.\end{propriete}
\begin{definition}Une grandeur produit se calcule en multipliant deux grandeurs. Par exemple, à puissance constante, l’énergie est $E=P\times t$. Avec $P$ en kW et $t$ en h, $E$ est en kWh. Les formules scientifiques utiles sont données dans les exercices.\end{definition}
\begin{savoirfaire}[Contrôler]Écrire les unités à côté des données avant le calcul. Vérifier si une augmentation doit faire augmenter ou diminuer le résultat : pour un même volume, un débit plus fort donne une durée plus courte. Donner un ordre de grandeur et une phrase avec l’unité finale.\end{savoirfaire}

\clearpage\rubrique{Découverte}
\begin{blocLivre}{ex-8-1}\begin{exercice}[1]{Une durée}
Convertir 2 h 18 min en heures décimales et en minutes.
\end{exercice}\end{blocLivre}
\begin{blocLivre}{ex-8-2}\begin{exercice}[1]{Une vitesse moyenne}
Un cycliste parcourt 45 km en 2 h 30 min. Calculer sa vitesse moyenne.
\end{exercice}\end{blocLivre}
\begin{blocLivre}{ex-8-3}\begin{exercice}[1]{Une conversion de vitesse}
Convertir 72 km/h en m/s.
\end{exercice}\end{blocLivre}
\begin{blocLivre}{ex-8-4}\begin{exercice}[1]{Un débit}
Un robinet délivre 18 L en 45 s. Donner son débit en L/s.
\end{exercice}\end{blocLivre}
\begin{blocLivre}{ex-8-5}\begin{exercice}[1]{Un volume et une contenance}
Convertir $0{,}35$ m$^3$ en litres et 250 mL en cm$^3$.
\end{exercice}\end{blocLivre}

\rubrique{Entraînement}
\begin{blocLivre}{ex-8-6}\begin{exercice}[2]{Distance et temps}
À vitesse constante de 12 km/h, quelle distance parcourt-on en 35 min ?
\end{exercice}\end{blocLivre}
\begin{blocLivre}{ex-8-7}\begin{exercice}[2]{Remplir une cuve}
Une cuve de 450 L est remplie à 15 L/min. Combien de temps faut-il ?
\end{exercice}\end{blocLivre}
\begin{blocLivre}{ex-8-8}\begin{exercice}[2]{Changer les deux unités}
Convertir $0{,}012$ m$^3$/s en L/min.
\end{exercice}\end{blocLivre}
\begin{blocLivre}{ex-8-9}\begin{exercice}[2]{Une concentration}
On dissout 18 g de soluté dans une solution de volume final 300 mL. Calculer la concentration en g/L.
\end{exercice}\end{blocLivre}
\begin{blocLivre}{ex-8-10}\begin{exercice}[2]{Une énergie}
Un appareil de 1,2 kW fonctionne pendant 45 min. Avec $E=P\times t$, calculer l’énergie en kWh.
\end{exercice}\end{blocLivre}

\rubrique{Problèmes}
\begin{blocLivre}{ex-8-11}\begin{exercice}[3]{Aller et retour}
Un trajet de 30 km est parcouru à 15 km/h à l’aller et à 30 km/h au retour. Calculer la vitesse moyenne sur l’ensemble.
\end{exercice}\end{blocLivre}
\begin{blocLivre}{ex-8-12}\begin{exercice}[3]{Une masse volumique}
Une pièce a un volume de 40 cm$^3$ et une masse de 108 g. Calculer sa masse volumique, puis la masse d’une pièce de 65 cm$^3$ du même matériau.
\end{exercice}\end{blocLivre}
\begin{blocLivre}{ex-8-13}\begin{exercice}[3]{Deux arrivées d’eau}
Deux robinets ont des débits constants de 8 et 12 L/min. Ils remplissent ensemble un bassin de 1,2 m$^3$, initialement vide, sans fuite. Donner le temps de remplissage.
\end{exercice}\end{blocLivre}
\begin{blocLivre}{ex-8-14}\begin{exercice}[3]{Une pause dans le trajet}
Un marcheur parcourt 9 km à 4,5 km/h et fait une pause de 30 min. Quelle est sa vitesse moyenne si l’on inclut la pause ?
\end{exercice}\end{blocLivre}
\begin{blocLivre}{ex-8-15}\begin{exercice}[3]{Une comparaison de débits}
La pompe A débite 900 L/min ; B débite $0{,}018$ m$^3$/s. Laquelle est la plus rapide ? Combien de minutes gagne-t-elle pour 5400 L ?
\end{exercice}\end{blocLivre}

\clearpage
